\documentclass[a4paper,11pt]{book}

%Main packages
\usepackage[T1]{fontenc}
\usepackage[utf8]{inputenc}
\usepackage{amsmath} %Math package.
\usepackage[bookmarks = true, hidelinks]{hyperref}
\usepackage{mathtools} %amsmath enhancement.
\usepackage{amsthm} %Theorem environments.
\usepackage{amssymb} %Extended symbol collection.
\usepackage{mathrsfs} %\mathscr.
\usepackage[backend=biber, style=alphabetic]{biblatex}
\addbibresource{bibliography.bib}
\usepackage[capitalise]{cleveref} %Defines \cref.
\usepackage[english]{babel} %Language.
\usepackage{enumitem} %Used for \itemize.
\usepackage{lipsum} %Lorem ipsum generator.
\usepackage{graphicx} %Figures.
\usepackage{tikz-cd} %Commutative diagrams.
\usetikzlibrary{calc, matrix} %Used from tikz package.
\usepackage{xcolor} %Colored text.
\usepackage{mathpazo}
\usepackage[scaled=0.95]{helvet}
\usepackage{courier}
\usepackage{eucal} %Used to replace the ugly calligraphic by something better.
\DeclareMathAlphabet{\mathpzc}{OT1}{pzc}{m}{it} %\mathpzc{} gives lowercase calligraphic symbols.
\usepackage{spectralsequences}

%Geometry
\usepackage[a4paper,
	left = 1in, 
	right = 1in,
	top = 1in,
	bottom = 1.3in,
	marginparwidth = 1.9in,
	marginparsep = 0.2in,
	headsep = 30pt]{geometry} %Redefine margins.
\setlength{\parindent}{20pt} %Indentation of new paragraph.
\linespread{1.15} %Spread out the lines a bit more.
\widowpenalty10000 %Prevent one-line paragraphs.
\clubpenalty10000 %Prevent one-line paragraphs.
\raggedbottom %Might fix some equation spacing.
\sloppy %Kills hbox overflows

%\setcounter{secnumdepth}{0} %Add if we don't want numbered sections

%Environments
%\numberwithin{equation}{section} %Uncomment this to get more detailed theorem numberings
\theoremstyle{definition}
\newtheorem{thm}[equation]{Theorem}
\newtheorem{lem}[equation]{Lemma}
\newtheorem{prop}[equation]{Proposition}
\newtheorem{conj}[equation]{Conjecture}
\newtheorem{ex}[equation]{Example}
\newtheorem{cor}[equation]{Corollary}
\newtheorem{rmk}[equation]{Remark}

%\numberwithin{figure}{section} %Uncomment these to get more detailed figure & table numberings
%\numberwithin{table}{section} %Uncomment these to get more detailed figure & table numberings
\makeatletter 
\let\c@table\c@figure %Match up figure and table numbering
\makeatother


%New commands --- here are some examples
\newcommand{\mc}{\mathcal}
\newcommand{\ra}{\rightarrow}
\newcommand{\s}{\subseteq}
\newcommand{\bb}{\mathbb}
\newcommand{\name}{\operatorname}
\newcommand{\ol}{\overline}
\newcommand{\Ra}{\Longrightarrow}
\DeclareMathOperator{\Hom}{Hom}
\newcommand{\operator}[1]{\begingroup\mathgroup-1 #1\endgroup} %Use this to declare operators with uncommon symbols
\DeclareMathOperator{\ho}{\operator{\pi}}
\DeclareMathOperator{\holim}{\operator{\varprojlim}}

\author{Name of the author\\
	\today}
\title{\Huge A fancy title}
\date{}

\begin{document}

\frontmatter
\pagestyle{plain}

\maketitle

\chapter{Abstract}

\lipsum[1-2]

\section*{Acknowledgements}

\lipsum[3]

\vspace{\fill}
\noindent
\framebox[\textwidth]{
\begin{tabular*}{0.96\textwidth}{@{\extracolsep{\fill} }cp{0.84\textwidth}}
 % The EU emblem
\raisebox{-0.7\height}{
\begin{tikzpicture}[y = 0.80pt, x = 0.8pt, yscale = -1, inner sep = 0pt, outer sep = 0pt, scale = 0.12]
    \definecolor{c003399}{RGB}{0,51,153}
    \definecolor{cffcc00}{RGB}{255,204,0}
    \begin{scope}[shift={(0,-872.36218)}]
        \path[shift={(0,872.36218)},fill=c003399,nonzero rule] (0.0000,0.0000) rectangle (270.0000,180.0000);
        \foreach \myshift in 
            {(0,812.36218), (0,932.36218), 
            (60.0,872.36218), (-60.0,872.36218), 
            (30.0,820.36218), (-30.0,820.36218),
            (30.0,924.36218), (-30.0,924.36218),
            (-52.0,842.36218), (52.0,842.36218), 
            (52.0,902.36218), (-52.0,902.36218)}
        \path[shift=\myshift,fill=cffcc00,nonzero rule] (135.0000,80.0000) -- (137.2453,86.9096) -- (144.5106,86.9098) -- (138.6330,91.1804) -- (140.8778,98.0902) -- (135.0000,93.8200) -- (129.1222,98.0902) -- (131.3670,91.1804) -- (125.4894,86.9098) -- (132.7547,86.9096) -- cycle;
    \end{scope}
\end{tikzpicture}}
&
\footnotesize{This project has received funding from the European Union Horizon 2020 research and innovation programme under the Marie Sk{\l}odowska--Curie grant agreement N\textsuperscript{\underline{o}} 801199.}
\end{tabular*}
}

\tableofcontents

\mainmatter

\chapter{First chapter}

\lipsum[3]

\section{A section}

\lipsum[4-5]

\section{Another section}

\lipsum[6-7]

\section{One more section}

\lipsum[8]

\begin{lem}
    \lipsum[1]
\end{lem}

\lipsum[9]


\begin{figure}
\centering
\begin{sseqpage}[
Adams grading,
yscale =0.8,
xscale =0.8,
x tick step = 2, 
x range = {-3}{5},
y range = {0}{8},
x label = {$t - s$},
y label = {$s$},
grid = chess, 
classes = { minimum width = 0.2cm, minimum height = 0.2cm },
run off = .-.]

\class[draw = none, "C_2"](0, 2)
\class[draw = none, "k^\times"](1, 0)
\class[diamond, fill](-1, 4)
\class(3, 0)
\class(5, 0)
\class["?"](-1, 2)
\class["?"](-1, 3)
\class["?"](-2, 4)
\class["?"](-2, 5)
\class["?"](-3, 4)
\class["?"](-3, 5)
\class["?"](-3, 6)

\class(-1, 6)
\class(-1, 8)
\class(-1, 6)
\class(-1, 8)
\class(-2, 7)
\class(-2, 7)
\class(-3, 8)
\class(-3, 8)
\class(0, 3)
\class(0, 5)
\class(0, 7)
\class(0, 3)
\class(0, 5)
\class(0, 7)
\class(1, 2)
\class(1, 4)
\class(1, 6)
\class(1, 8)
\class(1, 2)
\class(1, 4)
\class(1, 6)
\class(1, 8)
\class(2, 1)
\class(2, 3)
\class(2, 5)
\class(2, 7)
\class(2, 1)
\class(2, 3)
\class(2, 5)
\class(2, 7)
\class(3, 2)
\class(3, 4)
\class(3, 6)
\class(3, 8)
\class(3, 2)
\class(3, 4)
\class(3, 6)
\class(3, 8)
\class(4, 1)
\class(4, 3)
\class(4, 5)
\class(4, 7)
\class(4, 1)
\class(4, 3)
\class(4, 5)
\class(4, 7)
\class(5, 2)
\class(5, 4)
\class(5, 6)
\class(5, 8)
\class(5, 2)
\class(5, 4)
\class(5, 6)
\class(5, 8)

\d[dashed]3(0, 3)
\d3(1, 2)
\d3(1, 2, 2, 2)
\d3(2, 1)
\d3(2, 1, 2, 2)
\d3(3, 0)
\d3(2, 5)
\d3(2, 5, 2, 2)
\d3(3, 4)
\d3(3, 4, 2, 2)
\d3(4, 3)
\d3(4, 3, 2, 2)
\d3(5, 2)
\d3(5, 2, 2, 2)

\class(2, 11)
\class(2, 11)
\class(3, 10)
\class(3, 10)
\class(4, 9)
\class(4, 9)
\class(6, 1)
\class(6, 1)
\class(6, 5)
\class(6, 5)
\d3(3, 8)
\d3(3, 8, 2, 2)
\d3(4, 7)
\d3(4, 7, 2, 2)
\d3(5, 6)
\d3(5, 6, 2, 2)
\d3(6, 1)
\d3(6, 1, 2, 2)
\d3(6, 5)
\d3(6, 5, 2, 2)
\class(0, 9)\class(0, 9)
\class(-1, 10)\class(-1, 10)
\class(-2, 11)\class(-2, 11)
\d3(1, 6)
\d3(1, 6, 2, 2)
\d3(0, 7)
\d3(0, 7, 2, 2)
\d3(-1, 8)
\d3(-1, 8, 2, 2)
\end{sseqpage}
    \caption{$E_3$-page of the limit spectral sequence for computing the Picard group of $\mathsf{StMod}(kQ_8)$. The nontrivial differentials have been illustrated.}
\label{E3limitQ8}
\end{figure}

\lipsum[10]

\begin{equation}
    \#X(\bb{F}_q) = \sum_{i=0}^n (-1)^i \name{Tr}{\big(\name{Frob}_q^* \mid H^i_c(\ol{X};\bb{Q}_\ell)\big)}\text{.}
\end{equation}

\lipsum[9]

\begin{thm}
    \lipsum[3]
\end{thm}

\begin{proof}
    \lipsum[2-5]
\end{proof}

\chapter{Second chapter}

\lipsum[11]

\begin{lem}[{\cite[Theorem 2.1]{riemann1919hypothesen}}]
    \lipsum[3]
\end{lem}

\section{A fourth section}

\lipsum[2-3]

\section{Yet another section}

\lipsum[4-7]

\section{The last section}

\lipsum[8]
\[ E_{2}^{st} = H^s(\mc{I};\ho_t \mc{F}) \Ra \ho_{t-s} \holim_\mc{I} \mc{F} \text{.}\]

And here's some more text. \lipsum[9]

\addcontentsline{toc}{chapter}{References}
\printbibliography

\end{document}
